\chapter{Conclusions}

Throughout the three previous chapters, we have been able to provide efficient
certifying algorithms to solve \textsc{Monopolarity} on $(P_5,\text{house})$-free
graphs depending on whether the graph is chordal, or contains an induced copy of
either $C_4$ or $C_5$. Whenever the graphs contains an induced copy of $C_4$ or
$C_5$, the algotithm requires for such a copy to be given as an input.
Nonetheless, if the algorithm changes its state to execute the procedures of
either \Cref{cap:c5} or \Cref{cap:no_c5}, it happens because it has already
found an induced cycle. This result is summarized in the next corollary as a
consequence of \Cref{cor:algorithm_chordal,cor:algorithm_c5,cor:algorithm_c4}.

\begin{corollary}
    There is a linear-time certifying algorithm that decides whether a
    $(P_5,\text{house})$-free graph is monopolar.
\end{corollary}

Recall that, except for the graphs containing an induced copy of $C_4$, we were
able to describe the complete structure of monopolar $(P_5,\text{house})$-free
graphs. Thus, a slightly stronger result can be stated for those graphs, as a
consequence of \Cref{cor:algorithm_chordal,cor:algorithm_c5}.

\begin{corollary}
    There is a linear-time certifying algorithm that decides if a $C_4$-free
    graph is monopolar and $(P_5,\text{house})$-free.
\end{corollary}

We leave as an open problem to complete the structure of monopolar $(P_5,
\text{house})$-free graphs with an induced $C_4$, which in turn will generalize
the last corollary to graphs that are not $C_4$-free. Once that is done, we
would have a linear-time certifying recognition algorithm for the family of
monopolar $(P_5,\text{house})$-free graphs. For the moment, observe that our
algorithm still can return a monopolar partition even when the input graph is
not $(P_5,\text{house})$-free. Whenever the graph is chordal, we have ensured
that it also returns an induced copy of $P_5$, if one exists, and the same goes
for graphs with an induced copy of $C_5$. Lastly, if the algorithm outputs one
of the monopolar partitions given in the proof of \Cref{thm:monopolar_c4}, then
we cannot ensure that the input graph is indeed $(P_5,\text{house})$-free.

Having characterized monopolarity in $(P_5,\text{house})$-free graphs, we can
now consider the more general problem of polarity in the same class of graphs.
In particular, the next result proved in \cite{Contreras2024} implies that we
already have all disconnected $(P_5,\text{house})$-free obstructions to
polarity.

\begin{lemma}[\cite{Contreras2024}]
    If $G$ is a graph, then $G$ is disconnected minimal polar obstruction if and
    only if $G \cong P_3 + H$, where $H$ is a minimal monopolar obstruction
    which is not a minimal polar obstruction.
\end{lemma}

It is easy to see that from all monopolar obstructions given, only $H_7$, $H_9$
and $H_{10}$ are not polar. Thus, the aforementioned lemma implies the following
result.

\begin{corollary}
    The disconnected $(P_5,\text{house})$-free minimal polar obstructions are
    the graphs $P_3 + H_i$ with $i\in \l{1,...,14}\setminus \l{7,9,10}$.
\end{corollary}

We leave as an open problem the study the connected $(P_5,\text{house})$-free
minimal polar obstructions.  Observe that it suffices to consider the connected
$(P_5,\text{house})$-free minimal polar obstructions that also have a connected
complement, so all obstructions left to find are not cographs (graphs with no
induced $P_4$).

We now briefley turn our attention to the complementary problem of
\textsc{Monopolarity}, that is, the problem of finding an $(\infty,1)$-polar
partition. Since the class of $(P_5,\text{house})$-free graphs is closed under
complementation, the following result holds as a direct consequence of
\Cref{cor:algorithm_chordal,cor:algorithm_c5,cor:algorithm_c4} and the fact that
the complement of a graph can be computed in time $O(|V|^2)$.

\begin{corollary}
    There is a $O(|V|^2)$ certifying algorithm that decides whether a
    $(P_5,\text{house})$-free graph admits an $(\infty,1)$-polar partition.
    Additionally, the minimal obstructions to $(\infty,1)$-polarity in the class
    of $(P_5,\text{house})$-free graphs are the graphs $\overline{H_i}$ with
    $i\in \l{1,...,14}$.
\end{corollary}

We believe that a linear-time certifying algorithm to recognize
$(\infty,1)$-polar graphs when the input is a $(P_5, \text{house})$-free graph
can be obtained by means of an analysis similar to the one done in this work.
Nonetheless, some details must be taken into consideration, in particular, the
tests that involve calculating connected components and their interaction.

Finally, we give an additional result regarding intermediate levels of polarity.
Let $k$ be a positive integer.   The decision problem
$(k,\infty)$-\textsc{Polarity} takes a graph $G$ as an input and asks whether
$G$ admits a $(k,\infty)$-polar partition.   Clearly
$(1,\infty)$-\textsc{Polarity} coincides with \textsc{Monopolarity}.

\begin{proposition}
\label{pro:np-hardness}
    The problem $(k,\infty)$-\textnormal{\scshape Polarity} is NP-complete for
    any positive integer $k$.
\end{proposition}

\begin{proof}
    We proceed by induction. Since \textsc{Monopolarity} is NP-complete, we have
    that $(1,\infty)$-\textsc{Polarity} is NP-complete.

    Suppose that $(k,\infty)$-\textsc{Polarity} is NP-complete. By means of a
    polynomial reduction, we proceed to prove that
    $(k+1,\infty)$-\textsc{Polarity} is also NP-complete. Clearly this last
    problem is NP, since a polynomial time verifier for this language receives
    the input graph $G$ along with the $(k+1,\infty)$-polar partition of $G$,
    $(A,B)$, if it exists. The process of verification searches for an induced
    copy of either $K_2 + K_1$ or $K_{k+2}$ in $G[A]$, and searches for an
    induced copy of $P_3$ in $G[B]$.  Both searches can be performed in time
    $O(n^{k+2})$.

    Let $G$ be an input graph for $(k,\infty)$-\textsc{Polarity}.  The
    polynomial reduction $G'$ is obtained from $G$ by adding two vertices, $x$
    and $y$, which are universal to $G$ and nonadjacent to each other.    It is
    immediate that $G'$ can be constructed in time $O(|V_G|)$. We claim that $G$
    has a $(k,\infty)$-polar partition if and only if $G'$ has a
    $(k+1,\infty)$-polar partition.

    Suppose that $G$ has a $(k,\infty)$-polar partition, $(A,B)$. Since $x$ and
    $y$ are nonadjacent and complete to $G$, it follows that $A\cup \l{x,y}$
    induces a complete $(k+1)$-partite graph. Hence, $(A\cup \l{x,y}, B)$ is a
    $(k+1,\infty)$-polar partition of $G'$.

    Suppose that $G'$ has a $(k+1,\infty)$-polar partition $(A,B)$. We proceed
    by cases depending on which part $x$ and $y$ lie on.

    \paragraph{Case 1: Both $x$ and $y$ are in $B$.} Since $x$ and $y$ are
    nonadjacent, complete to $G$ and $B$ induces a cluster, we have that no
    other vertex of $G$ is in $B$. Therefore, $G=G'[A]$, so it is a complete
    $(k+1)$-partite graph. Let $A_1$ be a part of such a graph. Then,
    $A\setminus A_1$ induces a complete $k$-partite graph, while $A_1$ being
    stable implies that it induces a cluster, so $(A\setminus A_1, A_1)$ is a
    $(k,\infty)$-polar partition of $G$.

    \paragraph{Case 2: At least one of $x$ and $y$ is in $B$.} Suppose without
    loss of generality that $y\in B$. Since $B$ induces a cluster and $y$ is
    complete to $G$, it follows that $B$ has exactly one clique, $C$, which
    contains $y$, and $C\setminus \l{y}$ is a subset of $V_G$. Hence, all
    vertices of $G$, except for those in $C\setminus \l{y}$, are in $A$. Since
    $x\in A$ and it is complete to $G$, it follows that $x$ is in a singleton
    part of the complete $(k+1)$-partite graph $G'[A]$. Then, $G - (C\setminus
    \l{y})$ is a complete $k$-partite graph, and $C\setminus \l{y}$, which
    contains the remaining vertices of $G$, is a clique. Thus, $(A\setminus
    \l{x}, B\setminus \l{y})$ is a $(k,\infty)$-polar partition of $G$.

    \paragraph{Case 3: Both $x$ and $y$ are in $A$.} Since $x$ and $y$ are
    nonadjacent, it must be that $x$ and $y$ are in the same stable part of
    $G'[A]$.   Moreover, as $x$ and $y$ are complete to $G$, it must be the case
    that $\{x,y\}$ is a part in the complete multipartition of $G'[A]$.
    Therefore, $G'[A] - \{x, y\}$ is a complete $k$-partite graph, so
    $(A\setminus \l{x,y},B)$ is a $(k,\infty)$-polar partition of $G$.

    Thereby, deciding wheter a graph has a $(k+1,\infty)$-polar partition is
    NP-complete. The result follows by induction.

\end{proof}

The construction described in \Cref{pro:np-hardness} is such that it preserves
the absence of certain structures within graphs. For instance, consider a graph
$F$ with no universal vertex nor a pair of nonadjacent universal vertices. It is
clear that any $F$-free graph $G$ is such that the $G'$ constructed in the proof
of the last result is also $F$-free. In particular, for any disconnected graph
$F$ with at least one non-trivial component, $G$ being $F$-free implies that
$G'$ is $F$-free. The same holds for all graphs $F$ with every vertex having at
least two non-neighbors, which includes all cycles $C_n$ for $n\geq 5$. This
discussion allows us to further generalize the last proposition to families with
forbidden graphs in which monopolarity is NP-complete and the forbidden graphs
are as described above.   The families considered in our next result are just a
representative sample not intended to be exhaustive.

\begin{corollary}
    The problem $(k,\infty)$-\textnormal{\scshape Polarity} is NP-complete for
    any positive integer $k$ in the class of $\mathcal{F}$-free graphs, where
    $\mathcal{F}$ is one of the following:
    \begin{enumerate}
        \item $\mathcal{F} = \l{C_5,\text{house}}$
        \item $\mathcal{F} = \l{C_6,\text{house}}$
        \item $\mathcal{F} = \l{C_7,\text{odd anti-hole}}$
        \item $\mathcal{F} = \l{K_2 + K_3,\overline{P},\text{house}}$
        \item $\mathcal{F} = \l{P}$
        \item $\mathcal{F} = \l{S_3,S_4,\text{net}}$
        \item $\mathcal{F} = \l{\text{bull},\text{house}}$
    \end{enumerate}
\end{corollary}